%! Sinusoidal Function Context and Data Modeling — Unit 3, Lesson 3.7 — Guided Notes
\documentclass[10pt]{article}
\usepackage{precalculus-article}
\usepackage{precalculus-boxes}
\newcommand{\TallMath}[1]{$\displaystyle #1\rule[-1.4em]{0pt}{3.2em}$}

\begin{document}

\pageheader{Unit 3, Lesson 3.7}{Guided Notes}
\namedateperiod

\vspace{0.10in}

%----------------------------------------------------------------------
% Objectives
%----------------------------------------------------------------------
\begin{objectivebox}
\textbf{I will be able to\ldots}
\begin{enumerate}[leftmargin=1.5em, itemsep=4pt]
  \item \textbf{LO 3.7.A} --- Construct a sinusoidal model from a data set by
    computing amplitude, midline, period, and phase shift, and validate the
    model against known data values. \writeline
  \item \textbf{LO 3.7.A} --- Use a sinusoidal model to predict output values
    and identify the contextual domain over which the model is meaningful.
    \writeline
\end{enumerate}
\end{objectivebox}

\vspace{0.08in}

%----------------------------------------------------------------------
% Vocabulary
%----------------------------------------------------------------------
\begin{vocabbox}
\begin{description}[leftmargin=0pt, itemsep=6pt]
  \item[\termblanklong{sinusoidal model}]
    A function of the form $f(x) = A\sin(B(x - C)) + D$ (or cosine form) used
    to represent a \underline{\hspace{3cm}} real-world quantity.

  \item[\termblanklong{amplitude}]
    $A = \dfrac{\max - \min}{2}$; measures how far the function deviates from
    its \underline{\hspace{2cm}}.

  \item[\termblanklong{vertical shift / midline}]
    $D = \dfrac{\max + \min}{2}$; the horizontal line $y = D$ about which the
    function \underline{\hspace{3cm}}.

  \item[\termblanklong{period}]
    The \underline{\hspace{3cm}} input-value interval between consecutive maxima
    (or minima) in the data. Equals $\dfrac{2\pi}{B}$.

  \item[\termblanklong{phase shift}]
    $C$; shifts the graph \underline{\hspace{1.5cm}} when $C > 0$. Determined
    by comparing a data maximum to the standard model maximum location.

  \item[\termblanklong{contextual domain}]
    The set of input values over which the model is \underline{\hspace{3cm}}
    given the real-world context.
\end{description}
\end{vocabbox}

\vspace{0.08in}

%----------------------------------------------------------------------
% Hook
%----------------------------------------------------------------------
\begin{hookbox}
\textbf{Entry Question:} The table below shows average monthly high
temperatures (${}^\circ$F) for a city over one year.

\vspace{4pt}
\begin{center}
\renewcommand{\arraystretch}{1.3}
\small
\begin{tabular}{|c|c|c|c|c|c|c|c|c|c|c|c|c|}
\hline
\textbf{Month} $x$ & 1 & 2 & 3 & 4 & 5 & 6 & 7 & 8 & 9 & 10 & 11 & 12 \\
\hline
\textbf{Temp} $T(x)$ (${}^\circ$F) & 38 & 42 & 54 & 65 & 74 & 82 & 86 & 84 & 75 & 63 & 49 & 40 \\
\hline
\end{tabular}
\end{center}

\vspace{6pt}
\textbf{Q1.} Which month has the maximum temperature? \underline{\hspace{2cm}}
\quad Which has the minimum? \underline{\hspace{2cm}}

\vspace{4pt}
\textbf{Q2.} What kind of function could model this data?
\writelines{2}
\end{hookbox}

\vspace{0.08in}

%----------------------------------------------------------------------
% LO 3.7.A — Four-Step Model Building Process
%----------------------------------------------------------------------
\begin{notesbox}{LO 3.7.A --- Building a Sinusoidal Model: Four Steps}

We use the model $T(x) = A\sin(B(x - C)) + D$ where:
\begin{itemize}[leftmargin=1.4em, itemsep=2pt]
  \item $A =$ amplitude, \quad $D =$ midline (vertical shift),
  \item $B = \dfrac{2\pi}{k}$ ($k =$ period), \quad $C =$ phase shift.
\end{itemize}

\vspace{8pt}
\textbf{Step 1 --- Find the Period $k$}

The period equals the input-value distance between consecutive \underline{\hspace{2.5cm}}
(or consecutive \underline{\hspace{2.5cm}}).

In the temperature table, consecutive maxima occur at months \underline{\hspace{1.2cm}}
and \underline{\hspace{1.2cm}} (next year), so $k \approx$ \underline{\hspace{1.5cm}}.

$B = \dfrac{2\pi}{k} = $ \underline{\hspace{3cm}}

\vspace{10pt}
\textbf{Step 2 --- Compute Amplitude and Midline}

$\max = $ \underline{\hspace{1.5cm}} \quad $\min = $ \underline{\hspace{1.5cm}}

$A = \dfrac{\max - \min}{2} = \dfrac{\underline{\hspace{1cm}} - \underline{\hspace{1cm}}}{2} =$ \underline{\hspace{2cm}}

$D = \dfrac{\max + \min}{2} = \dfrac{\underline{\hspace{1cm}} + \underline{\hspace{1cm}}}{2} =$ \underline{\hspace{2cm}}

The midline represents \underline{\hspace{7cm}} in context.

\vspace{10pt}
\textbf{Step 3 --- Determine the Phase Shift $C$}

For $T(x) = A\sin(B(x - C)) + D$, the maximum occurs when $B(x - C) = \dfrac{\pi}{2}$,
so $x = C + \dfrac{\pi}{2B}$.

The data maximum occurs at $x = $ \underline{\hspace{1.2cm}}. Set equal and solve for $C$:

\vspace{4pt}
$C + \dfrac{\pi}{2B} = $ \underline{\hspace{1.5cm}} \quad $\Rightarrow$ \quad $C \approx$ \underline{\hspace{2cm}}

\vspace{6pt}
\textit{Tip:} If you use a \textbf{cosine} model $T(x) = A\cos(B(x-C)) + D$, the maximum
occurs when $B(x - C) = 0$, so $C = x_{\max}$ \underline{\hspace{3cm}}.

\vspace{10pt}
\textbf{Step 4 --- Write and Validate the Model}

$T(x) = $ \underline{\hspace{6cm}}

\vspace{6pt}
Validate: substitute a known data point (e.g., $x = 1$, $T = 38$):

$T(1) \approx$ \underline{\hspace{3cm}} \quad Residual $= $ actual $-$ predicted $=$ \underline{\hspace{2cm}}

\vspace{6pt}
\textbf{Contextual Domain:} This model is valid for \underline{\hspace{5cm}}.

\end{notesbox}

\vspace{0.08in}

%----------------------------------------------------------------------
% Practice
%----------------------------------------------------------------------
\begin{practicebox}

\textbf{Guided Practice.}\quad The table below shows the depth of water (in feet)
at a harbor entrance over one tidal cycle (12 hours).

\vspace{4pt}
\begin{center}
\renewcommand{\arraystretch}{1.35}
\begin{tabular}{|c|c|c|c|c|c|c|c|}
\hline
Time $t$ (hr) & 0 & 2 & 4 & 6 & 8 & 10 & 12 \\
\hline
Depth $d(t)$ (ft) & 5 & 9.5 & 13 & 9.5 & 5 & 0.5 & 5 \\
\hline
\end{tabular}
\end{center}

\vspace{4pt}
(a)\quad Identify $\max = $ \underline{\hspace{1.5cm}} at $t = $ \underline{\hspace{1cm}}
and $\min = $ \underline{\hspace{1.5cm}} at $t = $ \underline{\hspace{1cm}}.

\vspace{4pt}
(b)\quad Compute: $A = $ \underline{\hspace{2.5cm}} \quad $D = $ \underline{\hspace{2.5cm}}

\vspace{4pt}
(c)\quad Estimate the period $k = $ \underline{\hspace{1.5cm}} hr. \quad Then $B = \dfrac{2\pi}{k} = $ \underline{\hspace{2.5cm}}.

\vspace{4pt}
(d)\quad Identify the phase shift $C$, using a cosine model with maximum at $t = 4$. $C = $ \underline{\hspace{1.5cm}}.

\vspace{4pt}
(e)\quad Write the model: $d(t) = $ \underline{\hspace{7cm}}

\vspace{4pt}
(f)\quad Use your model to predict the depth at $t = 3$. \writeline

\vspace{4pt}
(g)\quad Is $t = 15$ hr inside the contextual domain for one tidal cycle? Explain.
\writelines{2}

\end{practicebox}

\end{document}
